\chapter{Dominios, codominios e hipótesis de los resultados}
\label{app:premisas}

\begin{enumerate}
\item El censo \(104\,976\to468\to243\), la imagen de 42 ternas, la
factorización \(468=26\times18\), las cinco regiones y los cinco ciclos son
teoremas exhaustivos del sistema finito de transporte especificado en esta
obra.

\item El catálogo enriquecido produce, por predicados finitos de fibra y sin
consulta decimal, la única órbita de clausura, la única órbita de
propagación y la única órbita de autoescala; el calibre interno orienta en
ellas los representantes \(010211,201101,121200\). Sus caracteres se
construyen, respectivamente, por la clausura pentarregional con retorno
unitario, la recurrencia semigrupal \((n+1)a_{n+1}=a_n\) y el rayo positivo
único de la incidencia \(F_{\rm av}\). La evaluación arquimediana publica
entonces los valores generados y el reconocimiento convencional los denomina
\(\pi,e,\varphi\). Para cada profundidad \(N\), la prolongación TPK produce
un prefijo único compatible con el de profundidad \(N-1\); el límite
inverso de esos prefijos fija la expansión completa. Por ello, las palabras,
los caracteres y los valores son etapas sucesivas de una construcción
interna, no etiquetas elegidas mediante tablas externas.

\item El teorema dodecafásico tiene por dominio causal el estado terminal
enriquecido \(x_{\rm term}\) producido por APP--TRIT--TPK. La carta firmada
publica desde él
\(x_{\rm term}\xrightarrow{R_{\rm sgn}}U_{12}^{\rm sgn}\), y la inversión
integral publica después
\(U_{12}^{\rm sgn}\xrightarrow{\mathcal H_{12}^{-1}}K\). En este dominio,
las correspondencias con \((D_3K,D_4K,Q)\), la representación de Hadamard y
la clausura dodecafásica son exactas y reversibles. Ni \(U\) ni \(K\) son
entradas generadoras. La traza reducida de una trayectoria posicional y el
sistema bivariado completo son objetos distintos, según se establece en el
\cref{chap:tpk-definicion}.

\item La igualdad con la cadena central 2018 es un reconocimiento externo.
La palabra matemática interna y su demostración son independientes de las
recomendaciones metrológicas posteriores.

\item La instancia finita de la doble proyección demuestra una bandera común,
\(W_{12}\) y \(M_{12}\). En la rama reticular se fijan el pegado de
\(A_2^{12}\), la cámara radial y el origen; el 3-vecino resultante se certifica
par, unimodular y sin raíces, y se identifica clásicamente con Leech. La
construcción de \(V^\natural\) y del Monstruo es una consecuencia clásica. La
flecha \(W_{12}\to W_{24}\) se realiza mediante el diseño residual de una
dodecada y la elevación única de sus hexadas, relativamente a una elección
\((D,\ell)\); las identificaciones se realizan mediante aplicaciones, no
mediante igualdades de cardinalidad. La compatibilidad proyectiva de estos
datos y su doble realización sobre el estado enriquecido común se demuestran
en los
\cref{p12-83-c17-s2:thm:lector-excepcional-global-rev6,p12-83-c17-s2:thm:doble-realizacion-global-rev6}.

\item El teorema de cilindros produce un único real para cada historia de su
dominio. La ley de prolongación cubre coherentemente cada compacto y el
agotamiento dirigido cubre la recta completa, como demuestran los
\cref{p12-83-c18-s1:thm:sobreyectividad-evaluacion-compacto,p12-83-c18-s1:prop:agotamiento-recta-real-u024,p12-83-c18-s2:hmtch:thm-cociente-real}.
Por tanto, la evaluación es sobreyectiva y su cociente arquimediano es
\(\mathbb R\); la cobertura es una conclusión del generador HMT, no una
premisa adicional.

\item El producto nativo está demostrado en el sector canónico de palabras.
La extensión monoidal al espacio completo de historias y el entrelazamiento
con las filtraciones de supervivencia se formulan como aplicaciones distintas
en el \cref{chap:producto-nativo}.
\end{enumerate}
